\documentclass[10pt,a3paper,landscape]{article}
\usepackage[margin=6mm]{geometry}
\usepackage[T1]{fontenc}
\usepackage[utf8]{inputenc}
\usepackage[french]{babel}
\usepackage{lmodern}
\usepackage{microtype}
\makeatletter\def\input@path{{../../../Style/}}\makeatother
\NeedsTeXFormat{LaTeX2e}
\ProvidesPackage{TLPosterCarteMentale}[2026/08/03 Posters et cartes mentales SI]

\RequirePackage{tikz}
\RequirePackage[most]{tcolorbox}
\RequirePackage{xcolor}
\RequirePackage{amsmath,amssymb}
\RequirePackage{graphicx}
\usetikzlibrary{positioning,arrows.meta,calc,shapes.geometric,mindmap,backgrounds,fit}

\definecolor{TLPosterBlue}{RGB}{24,86,141}
\definecolor{TLPosterLight}{RGB}{234,243,250}
\definecolor{TLPosterAccent}{RGB}{232,126,34}
\definecolor{TLPosterGreen}{RGB}{52,152,102}
\definecolor{TLPosterRed}{RGB}{192,57,43}
\definecolor{TLPosterGray}{RGB}{80,90,100}

\newcommand{\TLPosterTitre}[2]{%
  \begin{tcolorbox}[enhanced,boxrule=0pt,arc=0pt,colback=TLPosterBlue,left=6mm,right=6mm,top=4mm,bottom=4mm]
    {\color{white}\bfseries\LARGE #1}\hfill{\color{white}\large #2}
  \end{tcolorbox}%
}

\newtcolorbox{TLPosterBloc}[2][]{
  enhanced,breakable,
  colback=TLPosterLight,
  colframe=TLPosterBlue,
  colbacktitle=TLPosterBlue,
  coltitle=white,
  fonttitle=\bfseries,
  title={#2},
  boxrule=0.8pt,
  arc=2mm,
  left=3mm,right=3mm,top=2mm,bottom=2mm,
  #1
}

\newtcolorbox{TLPosterFormule}[1][]{
  enhanced,
  colback=white,
  colframe=TLPosterAccent,
  boxrule=1pt,
  arc=2mm,
  fontupper=\bfseries,
  halign=center,
  #1
}

\newcommand{\TLPosterMethode}[1]{%
  \begin{tcolorbox}[enhanced,colback=TLPosterGreen!8,colframe=TLPosterGreen,title=Méthode,fonttitle=\bfseries]
  #1
  \end{tcolorbox}%
}

\newcommand{\TLPosterAttention}[1]{%
  \begin{tcolorbox}[enhanced,colback=TLPosterRed!6,colframe=TLPosterRed,title=Point de vigilance,fonttitle=\bfseries]
  #1
  \end{tcolorbox}%
}

\tikzset{
  TLmindroot/.style={concept,concept color=TLPosterBlue,text=white,font=\bfseries\large,minimum size=34mm},
  TLmindlevel1/.style={level 1 concept/.append style={font=\bfseries,minimum size=23mm,sibling angle=45}},
  TLmindlevel2/.style={level 2 concept/.append style={font=\small,minimum size=17mm}},
  TLarrow/.style={-{Stealth[length=3mm]},thick,draw=TLPosterBlue},
  TLbox/.style={rounded corners=2mm,draw=TLPosterBlue,fill=TLPosterLight,align=center,inner sep=3mm}
}

\newenvironment{TLCarteMentale}[1]{%
  \begin{tikzpicture}[mindmap,grow cyclic,concept color=TLPosterBlue,every node/.style={concept},TLmindlevel1,TLmindlevel2]
  \node[TLmindroot]{#1}
}{;
  \end{tikzpicture}
}

\newcommand{\TLBranche}[3][]{child[concept color=#2]{node[#1]{#3}}}

\endinput

\pagestyle{empty}
\renewcommand{\familydefault}{\sfdefault}
\begin{document}
\begin{tikzpicture}
\TLMapBase{Lois entrée-sortie et transmetteurs}{Boucles $\rightarrow$ Loi en position $\rightarrow$ Loi en vitesse $\rightarrow$ Transmission}
\TLMapBranchTL{Savoirs}{\textbullet\ distinguer structures cinématiques ouvertes et fermées\\[1.2mm]\textbullet\ définir coordonnées articulaires et coordonnées opérationnelles\\[1.2mm]\textbullet\ établir une fermeture géométrique et une loi entrée-sortie en position\\[1.2mm]\textbullet\ établir une fermeture cinématique, dériver et projeter les relations\\[1.2mm]\textbullet\ connaître transmetteurs, rapports, rendement, couple, puissance et formule de Willis}
\TLMapBranchTR{Structures}{\textbullet\ ouverte / fermée\\[1.2mm]\textbullet\ coordonnées articulaires\\[1.2mm]\textbullet\ coordonnées opérationnelles}
\TLMapBranchML{Fermeture géométrique}{\textbullet\ boucle vectorielle\\[1.2mm]\textbullet\ projection\\[1.2mm]\textbullet\ loi en position}
\TLMapBranchMR{Vitesse}{\textbullet\ dérivation\\[1.2mm]\textbullet\ torseurs\\[1.2mm]\textbullet\ loi en vitesse}
\TLMapBranchBL{Transmetteurs et rapports}{\textbullet\ friction / courroies / chaînes\\[1.2mm]\textbullet\ engrenages / vis--écrou\\[1.2mm]\textbullet\ Willis / rendement / puissance}
\TLMapBranchBR{Méthode}{1. Choisir les paramètres d'entrée et de sortie et poser les repères.\\[1.2mm]2. Écrire une fermeture géométrique ou une contrainte adaptée.\\[1.2mm]3. Projeter et résoudre la loi en position.\\[1.2mm]4. Dériver ou composer les torseurs pour obtenir la loi en vitesse.\\[1.2mm]5. Établir le rapport de transmission avec la convention de signe.}
\TLMapRibbon{Modéliser $\rightarrow$ Écrire $\rightarrow$ Projeter $\rightarrow$ Résoudre $\rightarrow$ Dériver $\rightarrow$ Interpréter}
\end{tikzpicture}
\end{document}
